\paragraph{Interpretare.} Pentru toate capacitățile, ordinea este aceeași: CLOCK $\gtrsim$ LRU $>$ FIFO. La
$s = 0{,}99$ și 10\% capacitate: CLOCK \R{sim/zipf/0.99/10/clock}\%, LRU \R{sim/zipf/0.99/10/lru}\%, FIFO
\R{sim/zipf/0.99/10/fifo}\%, față de limita teoretică \R{sim/zipf/0.99/10/ideal}\%. Faptul că CLOCK egalează sau
depășește LRU este un rezultat cunoscut. Sub modelul IRM recența nu este semnalul optim, iar „a doua șansă"
acordată intrărilor referențiate acționează ca un filtru slab de frecvență. CLOCK este în plus mai ieftin,
pentru că un hit nu modifică lista.

Scanările secvențiale (rândul de jos) degradează toate politicile: la $s=0{,}99$ și 10\% capacitate, LRU
scade de la \R{sim/zipf/0.99/10/lru}\% la \R{sim/zipf+scan/0.99/10/lru}\%. Cheile reci, accesate o singură dată,
împing din cache setul fierbinte. Diferența de 5–14 puncte față de optim și sensibilitatea la scanări motivează
politicile cu admitere bazată pe frecvență (W-TinyLFU, ARC).
